
\documentclass[11pt,spanish]{article}

% =========================================================
% PLANTILLA BASE PARA INFORMES ACADÉMICOS / TÉCNICOS
% =========================================================

% --- Idioma y codificación ---
\usepackage[spanish,es-tabla]{babel}
\usepackage[T1]{fontenc}
\usepackage{textcomp}
\usepackage{newunicodechar}
\newunicodechar{₡}{\textcolonmonetary}
\newunicodechar{•}{\textbullet}
\newunicodechar{×}{\ensuremath{\times}}
\newunicodechar{–}{--}

% --- Márgenes / papel ---
\usepackage[
    left=2.5cm,
    right=2.5cm,
    top=2.5cm,
    bottom=2.5cm,
    headsep=0.5cm,
    footskip=0.8cm
]{geometry}

% --- Matemáticas ---
\usepackage{amsmath}

% --- Tablas, enlaces y elementos visuales ---
\usepackage{array}
\usepackage{tabularx}
\usepackage{booktabs}
\usepackage{longtable}
\usepackage{graphicx}
\usepackage{float}
\usepackage{xcolor}
\usepackage{colortbl}
\definecolor{tecblue}{HTML}{0B2F5B}
\definecolor{teclightblue}{HTML}{EAF1F8}
\definecolor{tablegray}{HTML}{F4F6F8}
\definecolor{diagramgreen}{HTML}{B2C19C}
\definecolor{diagramlight}{HTML}{E5E9DE}
\definecolor{diagramtext}{HTML}{565A60}
\usepackage{tikz}
\usetikzlibrary{positioning,shapes.geometric,arrows.meta,calc}
\usepackage{enumitem}
\usepackage{ragged2e}
\usepackage[hidelinks]{hyperref}

% --- Encabezado y pie de página ---
\usepackage{fancyhdr}
\usepackage{lastpage}

% --- Interlineado ---
\linespread{1.2}
\sloppy

% =========================================================
% DATOS DEL DOCUMENTO
% =========================================================
\newcommand{\Institucion}{Instituto Tecnológico de Costa Rica}
\newcommand{\Campus}{Campus Tecnológico Central Cartago}
\newcommand{\DocumentoTitulo}{Estimación de Costos}
\newcommand{\DocumentoSubtitulo}{Plataforma Universitaria Integrada: Modernización y Centralización de Procesos Académicos y Administrativos}
\newcommand{\Curso}{Administración de proyectos}
\newcommand{\Profesor}{Alicia Marcela Salazar Hernández}
\newcommand{\FechaDocumento}{16 de junio de 2026}
\newcommand{\PeriodoAcademico}{I Semestre}
\newcommand{\AnioDocumento}{2026}

\newcommand{\AutoresPortada}{%
    \begin{tabular}{l@{\hspace{1.5cm}}r}
        Mauricio González Prendas & 2024143009\\
        Isaac Jiménez Blanco & 2024202804\\
        Tamara Robles Camacho & 2024099342
    \end{tabular}%
}

% Datos que aparecen en el encabezado.
\newcommand{\DocHeaderLeft}{}
\newcommand{\DocHeaderRight}{Estimación de Costos}
\newcommand{\DocFooterRight}{Página~\thepage\ de~\pageref{LastPage}}

% Metadatos PDF.
\title{\DocumentoTitulo}
\author{\AutoresPortada}
\date{\FechaDocumento}

% =========================================================
% CONFIGURACIÓN DE TABLAS Y UTILIDADES
% =========================================================
\newcolumntype{Y}{>{\RaggedRight\arraybackslash}X}
\newcolumntype{P}[1]{>{\RaggedRight\arraybackslash}p{#1}}
\newcolumntype{R}[1]{>{\raggedleft\arraybackslash}p{#1}}
\renewcommand{\arraystretch}{1.22}
\setlength{\LTpre}{0.5em}
\setlength{\LTpost}{0.5em}

% Imagen segura: si el archivo no existe, muestra un recuadro de marcador.
\newcommand{\SafeImage}[2][0.92\textwidth]{%
    \IfFileExists{#2}{%
        \includegraphics[width=#1]{#2}%
    }{%
        \fbox{%
            \begin{minipage}[c][4cm][c]{#1}
            \centering
            Imagen pendiente:\\
            \texttt{#2}
            \end{minipage}%
        }%
    }%
}

% Figura reutilizable.
\newcommand{\EvidenceFigure}[3]{%
    \begin{figure}[H]
    \centering
    \SafeImage{#1}
    \caption{#2}
    \label{#3}
    \end{figure}%
}

% =========================================================
% ENCABEZADO Y PIE
% =========================================================
\pagestyle{fancy}
\fancyhf{}
\fancyhead[R]{\DocHeaderRight}
\renewcommand{\headrulewidth}{0.4pt}
\renewcommand{\footrulewidth}{0.4pt}
\fancyfoot[R]{\DocFooterRight}

% Estilo equivalente para páginas horizontales.
\fancypagestyle{widefancy}{%
    \fancyhf{}%
    \fancyhead[R]{\DocHeaderRight}%
    \renewcommand{\headrulewidth}{0.4pt}%
    \renewcommand{\footrulewidth}{0.4pt}%
    \fancyfoot[R]{\DocFooterRight}%
}

% =========================================================
% PÁGINAS HORIZONTALES PARA TABLAS ANCHAS
% =========================================================
\makeatletter
\newcommand{\SyncPageBuilderDimensions}{%
    \global\vsize=\textheight
    \global\@colroom=\textheight
    \global\@colht=\textheight
}
\makeatother

\newenvironment{widepage}{%
    \clearpage
    \hypersetup{pageanchor=false}%
    \paperwidth=297mm
    \paperheight=210mm
    \pdfpagewidth=297mm
    \pdfpageheight=210mm
    \newgeometry{left=2.5cm,right=2.5cm,top=2.5cm,bottom=2.5cm,headsep=0.5cm,footskip=0.8cm}%
    \setlength{\textwidth}{243mm}%
    \setlength{\textheight}{156mm}%
    \setlength{\linewidth}{\textwidth}%
    \setlength{\columnwidth}{\textwidth}%
    \hsize=\textwidth
    \setlength{\headwidth}{\textwidth}%
    \SyncPageBuilderDimensions
    \pagestyle{widefancy}%
}{%
    \clearpage
    \paperwidth=210mm
    \paperheight=297mm
    \pdfpagewidth=210mm
    \pdfpageheight=297mm
    \restoregeometry
    \SyncPageBuilderDimensions
    \hypersetup{pageanchor=true}%
    \pagestyle{fancy}%
}

% =========================================================
% PORTADA
% =========================================================
\newcommand{\MakeCover}{%
\begin{titlepage}
\thispagestyle{empty}
\centering

{\bfseries \huge \Institucion \par}
\vspace{0.25cm}
{\LARGE \Campus \par}

\vfill

{\huge \DocumentoTitulo \par}

\vfill

{\bfseries \LARGE Profesora: \par}
{\Large \Profesor \par}

\vfill

{\bfseries \LARGE Estudiantes: \par}
{\Large \AutoresPortada \par}

\vfill

{\bfseries\LARGE Fecha: \par}
{\Large \FechaDocumento \par}

\vfill

{\LARGE \PeriodoAcademico \par}
{\LARGE \AnioDocumento \par}

\end{titlepage}%
}

% =========================================================
% DOCUMENTO
% =========================================================
\begin{document}
\hypersetup{pageanchor=false}

% =========================
% Portada
% =========================
\MakeCover

% =========================
% Índice
% =========================
\pagenumbering{roman}
\pagestyle{empty}
\addtocontents{toc}{\protect\thispagestyle{empty}}
\tableofcontents
\clearpage
\pagenumbering{arabic}
\hypersetup{pageanchor=true}
\pagestyle{fancy}

% =========================
% Contenido
% =========================
\section{Información general del documento}

\begin{table}[H]
\centering
\begin{tabularx}{\textwidth}{p{5cm} Y}
\toprule
\textbf{Nombre del proyecto} & Plataforma Universitaria Integrada \\
\textbf{Organización} & Universidad Tecnológica La Mejor \\
\textbf{Documento} & \DocumentoTitulo \\
\textbf{Project manager} & Tamara Robles Camacho \\
\textbf{Moneda} & Colones costarricenses (₡) \\
\textbf{Técnicas aplicadas} & Bottom-Up (ascendente) y PERT (tres puntos) \\
\textbf{Versión} & 1.0 \\
\textbf{Fecha} & \FechaDocumento \\
\bottomrule
\end{tabularx}
\end{table}

\section{Introducción y metodología}

La estimación de costos es el proceso de predecir los recursos financieros necesarios para completar cada componente del proyecto. Para este proyecto se aplica una técnica (PERT) aprobada por el PMBOK:

\begin{table}[H]
\centering
\caption{Técnica aplicada}
\label{tab:tecnica-aplicada}
{\small
\begin{tabularx}{\textwidth}{P{3.8cm} Y Y}
\toprule
\textbf{Técnica} & \textbf{Descripción} & \textbf{Aplicación en este proyecto} \\
\midrule
PERT (Tres Puntos) & Considera incertidumbre con tres escenarios: Optimista (O), Más Probable (M) y Pesimista (P). Fórmula: (O + 4M + P) / 6 & Se aplica a los requisitos por semana, con sus tareas, además de los rubros adicionales: software, hardware, servicios y contingencias. \\
\bottomrule
\end{tabularx}
}
\end{table}

\section{Estimación PERT}

En la siguiente tabla se muestra la estimación de costos detallada para cada tarea utilizando la técnica PERT:

\begin{widepage}
{\footnotesize
\setlength{\tabcolsep}{3pt}
\begin{longtable}{P{7.9cm} R{3.0cm} R{3.2cm} R{3.0cm} R{4.8cm}}
\caption{Estimación PERT del proyecto}\label{tab:estimacion-pert}\\[-0.5em]
\toprule
\rowcolor{tecblue}
\multicolumn{5}{c}{\textcolor{white}{\textbf{Estimación PERT – Plataforma Universitaria Integrada}}} \\
\rowcolor{teclightblue}
\multicolumn{5}{c}{Universidad Tecnológica La Mejor  |  I Semestre 2026  |  Plazo: 3 meses  |  Equipo: 5 personas  |  Moneda: Colones (₡)} \\
\midrule
\rowcolor{tablegray}
\textbf{Tareas} & \textbf{Optimista} & \textbf{Más probable} & \textbf{Pesimista} & \textbf{Costo estimado} \par \textbf{(O + 4M + P) / 6} \\
\midrule
\endfirsthead
\toprule
\rowcolor{tablegray}
\textbf{Tareas} & \textbf{Optimista} & \textbf{Más probable} & \textbf{Pesimista} & \textbf{Costo estimado} \par \textbf{(O + 4M + P) / 6} \\
\midrule
\endhead
\rowcolor{diagramlight}
\multicolumn{5}{l}{\textbf{Fase 1: Inicio (Semanas 1–2)}} \\
\multicolumn{5}{l}{\textbf{\textit{Gestión de Inicio (Business Case, Charter, Stakeholders)}}} \\
\hspace{0.5cm}Elaborar Business Case completo & ₡480.000 & ₡700.000 & ₡900.000 & ₡696.667 \\
\hspace{0.5cm}Desarrollar Project Charter & ₡300.000 & ₡450.000 & ₡600.000 & ₡450.000 \\
\hspace{0.5cm}Registrar y analizar stakeholders & ₡200.000 & ₡280.000 & ₡420.000 & ₡290.000 \\
\hspace{0.5cm}Coordinación y gestión del PM (Inicio) & ₡300.000 & ₡450.000 & ₡600.000 & ₡450.000 \\
\midrule
\textbf{Subtotal} &  &  &  & \textbf{₡1.886.667} \\
\addlinespace[0.35em]
\rowcolor{diagramlight}
\multicolumn{5}{l}{\textbf{Fase 2: Planificación (Semanas 2–5)}} \\
\multicolumn{5}{l}{\textbf{\textit{Documentación PMI – Alcance, WBS, Cronograma}}} \\
\hspace{0.5cm}Definir alcance incluido/excluido & ₡200.000 & ₡280.000 & ₡420.000 & ₡290.000 \\
\hspace{0.5cm}Elaborar WBS con criterios de aceptación & ₡280.000 & ₡360.000 & ₡540.000 & ₡376.667 \\
\hspace{0.5cm}Cronograma en MS Project (tareas, hitos) & ₡350.000 & ₡480.000 & ₡720.000 & ₡498.333 \\
\multicolumn{5}{l}{\textbf{\textit{Documentación PMI – Recursos, Costos, Calidad}}} \\
\hspace{0.5cm}RACI, organigrama y Plan de Recursos & ₡200.000 & ₡280.000 & ₡420.000 & ₡290.000 \\
\hspace{0.5cm}Estimación de costos (PERT) & ₡180.000 & ₡240.000 & ₡360.000 & ₡250.000 \\
\hspace{0.5cm}Plan de Recursos & ₡110.000 & ₡220.000 & ₡330.000 & ₡220.000 \\
\hspace{0.5cm}Plan de Calidad con métricas & ₡160.000 & ₡220.000 & ₡330.000 & ₡228.333 \\
\multicolumn{5}{l}{\textbf{\textit{Documentación PMI – Comunicaciones, Riesgos, Adquisiciones}}} \\
\hspace{0.5cm}Plan de Comunicaciones & ₡140.000 & ₡200.000 & ₡300.000 & ₡206.667 \\
\hspace{0.5cm}Plan Gestión de Riesgos + Registro & ₡280.000 & ₡380.000 & ₡570.000 & ₡395.000 \\
\hspace{0.5cm}Matriz de Probabilidad e Impacto & ₡120.000 & ₡180.000 & ₡270.000 & ₡185.000 \\
\hspace{0.5cm}Plan de Adquisiciones & ₡120.000 & ₡180.000 & ₡270.000 & ₡185.000 \\
\hspace{0.5cm}Gestión y coordinación del PM (Planif.) & ₡400.000 & ₡540.000 & ₡810.000 & ₡561.667 \\
\midrule
\textbf{Subtotal} &  &  &  & \textbf{₡3.686.667} \\
\addlinespace[0.35em]
\rowcolor{diagramlight}
\multicolumn{5}{l}{\textbf{Fase 3: Ejecución – Desarrollo Frontend (Semanas 4–10)}} \\
\multicolumn{5}{l}{\textbf{\textit{Módulos del Sistema – Dev Frontend 1 (Mauricio)}}} \\
\hspace{0.5cm}Login y Dashboard Principal & ₡580.000 & ₡640.000 & ₡980.000 & ₡686.667 \\
\hspace{0.5cm}Portal Estudiantil – Matrícula y Horarios & ₡670.000 & ₡890.000 & ₡1.280.000 & ₡918.333 \\
\hspace{0.5cm}Portal Estudiantil – Calificaciones e Historial & ₡490.000 & ₡660.000 & ₡980.000 & ₡685.000 \\
\hspace{0.5cm}Portal Docente – Gestión cursos y asistencia & ₡500.000 & ₡700.000 & ₡960.000 & ₡710.000 \\
\hspace{0.5cm}Portal Docente – Registro de calificaciones & ₡360.000 & ₡490.000 & ₡720.000 & ₡506.667 \\
\hspace{0.5cm}Portal Administrativo – CRUD completo & ₡690.000 & ₡890.000 & ₡1.270.000 & ₡920.000 \\
\multicolumn{5}{l}{\textbf{\textit{Módulos del Sistema – Dev Frontend 2 (Carlos)}}} \\
\hspace{0.5cm}Portal Financiero – Pagos e historial & ₡480.000 & ₡700.000 & ₡960.000 & ₡706.667 \\
\hspace{0.5cm}Dashboard Ejecutivo del sistema & ₡480.000 & ₡640.000 & ₡960.000 & ₡666.667 \\
\multicolumn{5}{l}{\textbf{\textit{Módulos Compartidos – Ambos Devs}}} \\
\hspace{0.5cm}Integración entre módulos y navegación & ₡320.000 & ₡500.000 & ₡660.000 & ₡496.667 \\
\hspace{0.5cm}Pruebas de usabilidad y QA (ambos devs) & ₡480.000 & ₡710.000 & ₡960.000 & ₡713.333 \\
\midrule
\textbf{Subtotal} &  &  &  & \textbf{₡7.010.000} \\
\addlinespace[0.35em]
\rowcolor{diagramlight}
\multicolumn{5}{l}{\textbf{Fase 4: Monitoreo y Control (Semanas 1–12, continuo)}} \\
\multicolumn{5}{l}{\textbf{\textit{Seguimiento y Control}}} \\
\hspace{0.5cm}Definir y monitorear KPIs del proyecto & ₡200.000 & ₡350.000 & ₡420.000 & ₡336.667 \\
\hspace{0.5cm}Control de cambios + Solicitudes \#1 y \#2 & ₡240.000 & ₡320.000 & ₡480.000 & ₡333.333 \\
\hspace{0.5cm}Evaluación de impacto de cambios & ₡160.000 & ₡270.000 & ₡330.000 & ₡261.667 \\
\hspace{0.5cm}Seguimiento cronograma y riesgos activos & ₡200.000 & ₡280.000 & ₡420.000 & ₡290.000 \\
\hspace{0.5cm}Dashboard Ejecutivo del proyecto & ₡200.000 & ₡350.000 & ₡420.000 & ₡336.667 \\
\hspace{0.5cm}Gestión PM (reuniones, reportes, coordinación) & ₡480.000 & ₡710.000 & ₡960.000 & ₡713.333 \\
\midrule
\textbf{Subtotal} &  &  &  & \textbf{₡2.271.667} \\
\addlinespace[0.35em]
\rowcolor{diagramlight}
\multicolumn{5}{l}{\textbf{Fase 5: Cierre (Semana 12)}} \\
\multicolumn{5}{l}{\textbf{\textit{Entregables de Cierre}}} \\
\hspace{0.5cm}Informe Ejecutivo final & ₡120.000 & ₡220.000 & ₡270.000 & ₡211.667 \\
\hspace{0.5cm}Acta de Cierre formal & ₡100.000 & ₡190.000 & ₡210.000 & ₡178.333 \\
\hspace{0.5cm}Reflexiones y Recomendaciones del equipo & ₡120.000 & ₡200.000 & ₡270.000 & ₡198.333 \\
\midrule
\textbf{Subtotal} &  &  &  & \textbf{₡588.333} \\
\addlinespace[0.35em]
\rowcolor{diagramlight}
\multicolumn{5}{l}{\textbf{Rubros Adicionales (Software, Hardware, Servicios)}} \\
\multicolumn{5}{l}{\textbf{\textit{Software}}} \\
\hspace{0.5cm}MS Project licencia comercial (3 meses) & ₡30.000 & ₡70.000 & ₡90.000 & ₡66.667 \\
\hspace{0.5cm}Office 365 – 5 usuarios (3 meses) & ₡90.000 & ₡130.000 & ₡180.000 & ₡131.667 \\
\hspace{0.5cm}Herramientas premium opcionales & ₡0 & ₡20.000 & ₡60.000 & ₡23.333 \\
\multicolumn{5}{l}{\textbf{\textit{Hardware}}} \\
\hspace{0.5cm}USB / almacenamiento externo (5 unidades) & ₡10.000 & ₡20.000 & ₡25.000 & ₡19.167 \\
\hspace{0.5cm}Periféricos complementarios & ₡5.000 & ₡15.000 & ₡20.000 & ₡14.167 \\
\multicolumn{5}{l}{\textbf{\textit{Servicios}}} \\
\hspace{0.5cm}Internet (5 personas x 3 meses x ₡25,000) & ₡225.000 & ₡370.000 & ₡450.000 & ₡359.167 \\
\hspace{0.5cm}Impresiones y material físico & ₡6.000 & ₡10.000 & ₡17.000 & ₡10.500 \\
\hspace{0.5cm}Transporte a reuniones presenciales & ₡20.000 & ₡42.000 & ₡60.000 & ₡41.333 \\*
\midrule
\textbf{Subtotal} &  &  &  & \textbf{₡666.000} \\*
\midrule
Subtotal antes de contingencias &  &  &  & \textbf{₡16.109.333} \\*
Reserva para contingencias (10 \%) &  &  &  & \textbf{₡1.515.867} \\*
\midrule
\rowcolor{teclightblue}
\multicolumn{4}{l}{\textbf{Presupuesto total del proyecto (PERT + contingencias)}} & \textbf{₡17.625.200} \\*
\addlinespace[0.35em]
\rowcolor{tablegray}
\multicolumn{5}{l}{\textbf{Notas:}} \\
\multicolumn{5}{p{22.74cm}}{\scriptsize • Fórmula PERT: Costo Estimado = (Optimista + 4 × Más Probable + Pesimista) / 6} \\
\multicolumn{5}{p{22.74cm}}{\scriptsize • El escenario Optimista asume todo sale bien; el Pesimista incluye retrasos, correcciones y sobrecostos.} \\
\multicolumn{5}{p{22.74cm}}{\scriptsize • Equipo: Tamara Robles (PM), Mauricio González (Dev Frontend 1), Dev Frontend 2 (Carlos Sainz Arce), Isaac Jiménez (Analista 1), Analista 2 (Jorge Abarca Quiros).} \\
\multicolumn{5}{p{22.74cm}}{\scriptsize • Plazo total del proyecto: 3 meses (12 semanas). Moneda: Colones Costarricenses (₡).} \\
\multicolumn{5}{p{22.74cm}}{\scriptsize • Los costos de mano de obra se basan en tarifas de mercado TI Costa Rica para perfiles junior-medio (₡8,000–₡10,000/hr).} \\
\bottomrule
\end{longtable}
}
\end{widepage}

\section{Reserva para contingencias}

Según el estándar PMI, toda estimación debe incluir una reserva para contingencias que cubra riesgos identificados e imprevistos menores. Para este proyecto se aplica el 10\% sobre el subtotal de todos los rubros, que es la práctica estándar en proyectos de TI de mediana escala.

\begin{table}[H]
\centering
\caption{Reserva para contingencias}
\label{tab:reserva-contingencias}
\begin{tabularx}{\textwidth}{Y P{4cm}}
\toprule
\textbf{Concepto} & \textbf{Monto} \\
\midrule
Subtotal – Fase 1 (Inicio) & ₡1 886 667 \\
Subtotal – Fase 2 (Planificación) & ₡3 686 667 \\
Subtotal – Fase 3 (Desarrollo Frontend) & ₡7 010 000 \\
Subtotal – Fase 4 (Monitoreo y Control) & ₡2 271 667 \\
Subtotal – Fase 5 (Cierre) & ₡588 333 \\
Subtotal – Rubros Adicionales (Software, Hardware y Servicios) & ₡666 000 \\
Subtotal antes de contingencias & ₡16 109 333 \\
Reserva para contingencias (10 \%) & ₡1 515 867 \\
Presupuesto total del proyecto & ₡17 625 200 \\
\bottomrule
\end{tabularx}
\end{table}

\section{Resumen ejecutivo del presupuesto}

La siguiente tabla consolida todos los rubros identificados y presenta la distribución porcentual del presupuesto total del proyecto:

\begin{table}[H]
\centering
\caption{Resumen ejecutivo del presupuesto}
\label{tab:resumen-presupuesto}
\begin{tabularx}{\textwidth}{Y P{3.2cm} P{2.8cm} P{3cm}}
\toprule
\textbf{Rubro} & \textbf{Monto estimado} & \textbf{\% del total} & \textbf{Técnica usada} \\
\midrule
Fase 1 (Inicio) & ₡1 886 667 & 10.70\% & PERT \\
Fase 2 (Planificación) & ₡3 686 667 & 20.92\% & PERT \\
Fase 3 (Desarrollo Frontend) & ₡7 010 000 & 39.77\% & PERT \\
Fase 4 (Monitoreo y Control) & ₡2 271 667 & 12.89\% & PERT \\
Fase 5 (Cierre) & ₡588 333 & 3.34\% & PERT \\
Subtotal – Rubros Adicionales (Software, Hardware y Servicios) & ₡666 000 & 3.34\% & PERT \\
Contingencias (10\%) & ₡1 515 867 & 8.60\% & \% sobre subtotal \\
Presupuesto total del proyecto & 100\% & ₡17 625 200 \\
\bottomrule
\end{tabularx}
\end{table}

\subsection{Observaciones clave}

\begin{itemize}
    \item La mano de obra representa el 87.62\% del presupuesto total, lo que es típico en proyectos de software donde el capital humano es el principal activo.
    \item El uso de herramientas gratuitas (GitHub, React, VS Code) redujo significativamente el rubro de software.
    \item Las contingencias cubren principalmente retrasos en entregables documentales y ajustes de alcance no previstos.
\end{itemize}

\end{document}
